\section{Розділ~\ref{sec:acet}. Додаємо \nt{ACET}}

Три дії ацетилхоліну виміряно на зрізах і в живому мозку, конфлікт запису і читання показано на моделі книги, а те, що \nt{ACET} слугує лінією дозволу запису і повідомляє очікувану невизначеність, --- гіпотеза з частковою опорою на дослід.

{\footnotesize
\setlength{\tabcolsep}{3pt}\renewcommand{\arraystretch}{1.15}
\begin{longtable}{>{\raggedright}p{0.5cm}>{\raggedright}p{4.0cm}>{\raggedright}p{5.6cm}>{\raggedright}p{2.3cm}>{\raggedright\arraybackslash}p{1.7cm}}
\toprule
№ & Твердження & Дослід і що виміряно & Джерело & Статус \\
\midrule\endfirsthead
\toprule
№ & Твердження & Дослід і що виміряно & Джерело & Статус \\
\midrule\endhead
1 & \nt{ACET}-нейронів мозку небагато, вони зібрані в кількох групах, а виходи розходяться по всій корі: широкомовний канал & Забарвлення антитілами до ферменту синтезу ацетилхоліну і ретроградна мітка в щура: кора, гіпокамп, мигдалеподібне тіло, нюхова цибулина і таламус отримують ацетилхолін переважно від шести компактних груп клітин базального переднього мозку і стовбура & \cite{Mesulam1983} & виміряно \\
2 & Рівень ацетилхоліну в корі високий у неспанні і низький у повільному сні & Мікродіаліз у корі і гіпокампі вільно рухливих котів із записом ЕЕГ: викид ацетилхоліну в спокійному неспанні на $\sim100\%$, в активному --- на $\sim175\%$ вищий, ніж у повільному сні; у швидкому сні в корі --- як у неспанні & \cite{Marrosu1995,Hasselmo1999} & виміряно \\
3 & Зміст сигналу задає рецептор: ацетилхолін має швидкі рецептори-канали і повільні, що запускають реакції в клітині (твердження~\ref{prop:receptor}) & Огляд вимірювань: дія ацетилхоліну на збудливість, передачу і пластичність залежить від місця викиду, підтипу рецептора (нікотинові --- іонні канали, мускаринові --- через G-білки) і клітини-отримувача & \cite{Picciotto2012} & виміряно \\
4 & Перша дія: послабити внутрішні зв'язки, майже не зачіпаючи входів ззовні (множник $1-a$ у~\eqref{eq:acetnet}) & Зрізи нюхової кори щура: за $100$\,мкМ агоністів ацетилхоліну потенціали внутрішніх зв'язків (шар Ib) падають більш ніж на $60\%$, зовнішнього входу (шар Ia) --- менш ніж на $15\%$, у тій самій клітині; пригнічення пресинаптичне. Зрізи гіпокампа (CA1): $89\%$ проти $40\%$ для внутрішнього і вхідного шляхів & \cite{HasselmoBower1992,HasselmoSchnell1994} & виміряно \\
5 & Друга дія: посилити вхід (множник $\frac{1+a}{2}$) & Зрізи неокортексу: нікотинові рецептори посилюють синапси входу з таламуса і не діють на внутрішньокіркові; мускаринові послаблюють обидва типи. Збудливість нейронів ацетилхолін підвищує (огляд) & \cite{Gil1997,Hasselmo2006} & виміряно \\
6 & Третя дія: полегшити зміну ваг (швидкість $\eta a$) & Щур: електрична стимуляція базального ядра (джерела ацетилхоліну для кори) в парі з тоном спричиняє в дорослої тварини сильну перебудову карти слухової кори під звук, що пред'являється & \cite{KilgardMerzenich1998,Hasselmo2006} & виміряно \\
7 & Правило Гебба для рідкісних образів у другому рівнянні~\eqref{eq:acetnet} дає асоціативну пам'ять, що дописує образ за частиною & Розрахунок ємності мережі з розрідженими образами і коваріаційним правилом: за малої частки активних нейронів $p$ мережа зберігає помітно більше образів, ніж при $p=1/2$ & \cite{Tsodyks1988} & теорія \\
8 & Запис за низького $a$ псує пам'ять: мережа записує згадане, а не побачене; при $a=0.1$ псування не слабше, ніж при $a=0$ & \texttt{models/\allowbreak acet\_memory.py}: $200$ нейронів, п'ять образів по $20$, новий образ ділить з одним із них $10$ нейронів, $10$ прогонів. При $a=0$ новий образ не запам'ятано, старий тягне $5.9$ чужого нейрона з $10$; при $a=0.1$ новий образ згадується ($0.97$), але обидва зливаються: новий тягне $7.3$ чужого нейрона, старий --- $7.9$; з $a=0.2$ --- менше за один & виведено в розділі & модель \\
9 & Твердження~\ref{prop:writeread}: немає рівня $a$, за якого однаково добрі запис і читання (рис.~\ref{fig:acetsim}) & Той самий скрипт, швидкість навчання $\eta a$: новий образ запам'ятовується за одне пред'явлення цілком при $a\ge0.9$, на $0.8$ при $a=0.75$, не запам'ятовується при $a\le0.5$. Читання за половиною образу: $1.0$ при $a\le0.2$, $0.96$ при $a=0.3$, $0.04$ при $a=0.4$ & виведено в розділі & модель \\
10 & У мозку один широкомовний провід --- \nt{ACET} --- перемикає запис і читання & Ідею взято з моделей, побудованих за вимірюваннями на зрізах. Найпряміший дослід --- у людини: блокатор мускаринових рецепторів скополамін погіршує заучування нових пар слів, найсильніше тих, що перетинаються зі старими, але не згадування вже вивчених пар. Прямого вимірювання двох режимів мережі немає & \cite{HasselmoAndersonBower1992,Atri2004} & гіпотеза \\
11 & Фільтр~\eqref{eq:kalman} оптимальний; вага входу і швидкість навчання --- одна величина $P/R$, в усталеному режимі $P=\sqrt{qR}$ і пам'ять $\sqrt{R/q}$ (твердження~\ref{prop:kalman}) & Досліду не потребує: оптимальність фільтра Калмана--Б'юсі --- класичний результат теорії оцінювання; усталений режим виведено в розділі & \cite{KalmanBucy1961} & теорія \\
12 & \nt{ACET} повідомляє мережі величину, подібну до $P$, --- очікувану невизначеність & Запропоновано в імовірнісній моделі уваги. Прямого вимірювання того, що рівень ацетилхоліну йде за невизначеністю оцінки, немає & \cite{YuDayan2005} & гіпотеза \\
13 & Частина \nt{ACET}-нейронів відповідає на винагороду і покарання за два десятки мілісекунд, тим сильніше, чим несподіваніше підкріплення & Миша, оптогенетично впізнані холінергічні нейрони базального переднього мозку: відповідь на винагороду і покарання через $18\pm3$\,мс, величина зростає з несподіваністю підкріплення, відповіді схожі в нейронів двох ядер & \cite{Hangya2015} & виміряно \\
14 & \nt{NORA} помічає неочікувану зміну світу, \nt{ACET} тримає мережу в режимі запису & Розподіл праці запропоновано в моделі. Непрямо: у людини діаметр зіниці, пов'язаний із \nt{NORA}, відстежує зміни і швидкість навчання. Прямої перевірки пари \nt{NORA}--\nt{ACET} немає & \cite{YuDayan2005,Nassar2012} & гіпотеза \\
15 & У повільному сні мережа в режимі читання програє записане вдень, і ці повтори переписують його в довготривалу пам'ять & Записи ансамблів нейронів гіпокампа щура: клітини, що працювали разом удень, частіше спрацьовують разом у наступному повільному сні і в тому самому порядку --- виміряно. Для перезапису: пригнічення брижів гіпокампа після навчання погіршує просторову пам'ять, а подовження спонтанних брижів її покращує; що причина --- низький рівень \nt{ACET}, не доведено & \cite{WilsonMcNaughton1994,Skaggs1996,Girardeau2009,FernandezRuiz2019,Hasselmo1999} & гіпотеза \\
\bottomrule
\end{longtable}}
